\begin{example}
\label{example-roots-universal-polynomial}
Consider the ring map
\begin{eqnarray*}
R = \mathbf{Z}[a_1, \ldots, a_n]
& \longrightarrow &
S = \mathbf{Z}[\alpha_1, \ldots, \alpha_n] \\
a_1 & \longmapsto &
\alpha_1 + \ldots + \alpha_n \\
\ldots & \ldots & \ldots \\
a_n & \longmapsto & \alpha_1 \ldots \alpha_n
\end{eqnarray*}
In other words this is the unique ring map of polynomial
rings as indicated such that
$$
x^n + a_1 x^{n - 1} + \ldots + a_n
=
\prod\nolimits_{i = 1}^n (x + \alpha_i)
$$
holds in $\mathbf{Z}[\alpha_i, x]$. Another way to say this
is that $a_i$ maps to the $i$th elementary symmetric function
in $\alpha_1, \ldots, \alpha_n$. By the usual theory of elementary
symmetric polynomials (details omitted) the ring $S$ is finite
free over $R$ with basis the elements
$\alpha_1^{e_1} \alpha_2^{e_2} \ldots \alpha_n^{e_n}$ with
$0 \leq e_i \leq n - i$. A fortiori, the fibre rings of
$R \to S$ are finite and hence have dimension $0$.
On the other hand, $S$ is generated by
$n$ elements over $R$ subject to $n$ equations. Hence $S$
is a relative global complete intersection over $R$.
Since the rank of $S$ over $R$ is positive, we also see
that $S$ is faithfully flat over $R$.
\end{example}